\documentclass{ctexart}
\usepackage{amsmath, amssymb}
\usepackage{geometry}
\geometry{a4paper, margin=2cm}

\title{\textbf{第1章 质点运动学} --- 例题}
\author{}
\date{}

\begin{document}
\maketitle

\begin{enumerate}

\item 如图所示，以速度 $v$ 用绳跨一定滑轮拉湖面上的船，已知绳初长 $l_0$，岸高 $h$。求船的运动方程。

\item 已知一质点运动方程 $\vec{r} = 2t\,\vec{i} + (2 - t^{2})\,\vec{j}$。求：
\begin{enumerate}
  \item $t = 1\,\text{s}$ 到 $t = 2\,\text{s}$ 质点的位移；
  \item $t = 2\,\text{s}$ 时 $\vec{v}$、$\vec{a}$；
  \item 轨迹方程。
\end{enumerate}

\item（书中例题 1.1, 1.4，p.6; p.15）一质点作匀速圆周运动，半径为 $r$，角速度为 $\omega$。求用直角坐标、位矢表示的质点运动学方程。

\item（书中例题 1.2, 1.6，p.7; p.17，重点）直杆 $AB$ 两端可以分别在两固定且相互垂直的直导线槽上滑动，已知杆的倾角 $\varphi = \omega t$ 随时间变化，其中 $\omega$ 为常量。求杆中 $M$ 点的运动轨迹方程和加速度。

\item 已知 $a = 100 - 4t^{2}$，且 $t = 0$ 时 $v = 0$，$x = 0$。求速度 $v$ 和运动学方程。

\item 已知质点作匀加速直线运动，加速度为 $a$，求质点的运动方程。

\item（书中例题 1.9，p.21）在地球表面附近，质点以初速度 $v_0$ 倾斜抛出。不计空气阻力、风力、地球自转等影响，加速度即为重力加速度 $g$。设质点在 $Oxy$ 铅垂平面内作无阻力抛体运动，$a_x = 0$，$a_y = -g$。当 $t = t_0$ 时，$v_{0x} = v_0\cos\alpha$，$v_{0y} = v_0\sin\alpha$，$x_0 = y_0 = 0$，$g$、$v_0$ 均为常量。试求质点速度沿 $x$、$y$ 轴的投影和质点的运动学方程。

\item 一汽车在半径 $R = 200\,\text{m}$ 的圆弧形公路上行驶，其运动学方程为 $s = 20t - 0.2t^{2}\;(\text{SI})$。求汽车在 $t = 1\,\text{s}$ 时的速度和加速度大小。

\item 将一根光滑的钢丝弯成一个竖直平面内的曲线，质点可沿钢丝向下滑动。已知质点运动的切向加速度为 $a_{\tau} = -g\sin\theta$，$g$ 为重力加速度，$\theta$ 为切向与水平方向的夹角，初始速度 $v_0$，初始位置为 $y_0$。求质点在钢丝上各处的运动速度。

\item 一质点在水平面内以顺时针方向沿半径为 $2\,\text{m}$ 的圆形轨道运动。此质点的角速度与运动时间的平方成正比，即 $\omega = kt^{2}$，$k$ 为待定常数。已知质点在 $2\,\text{s}$ 末的线速度为 $32\,\text{m/s}$。求 $t = 0.5\,\text{s}$ 时质点的线速度和加速度。

\item（例 1.17）一细杆可绕过其一端的水平轴在铅直平面内自由转动。当杆与水平线夹角为 $\theta$ 时，其角加速度 $\beta = \dfrac{3g}{2l}\cos\theta$，$l$ 为杆长。试求：
\begin{enumerate}
  \item 杆自静止由 $\theta_0 = \dfrac{\pi}{6}$ 开始转至 $\theta = \dfrac{\pi}{3}$ 时杆的角速度；
  \item 杆的端点 $A$ 和中点 $B$ 的线速度大小。
\end{enumerate}

\item 一个带篷子的卡车，篷高 $h = 2\,\text{m}$，当它停在马路边时，雨滴可落入车内达 $d = 1\,\text{m}$，而当它以 $15\,\text{km/h}$ 的速率运动时，雨滴恰好不能落入车中。求雨滴的速度矢量。

\end{enumerate}

\end{document}
